\section*{Chapter \ref{ch:g4:volcanoes-earthquakes-first-look} --- Volcanoes and Earthquakes: A First Look}

\begin{solution}{exo:g4:volcanoes-earthquakes-first-look:1}
(a) \omterm{def:g4:volcanoes-earthquakes-first-look:eruption}{Lava}. (b) \omterm{def:g4:volcanoes-earthquakes-first-look:eruption}{Volcanic ash}. (c) Gas.
\end{solution}

\begin{solution}{exo:g4:volcanoes-earthquakes-first-look:2}
An \omterm{def:g4:volcanoes-earthquakes-first-look:earthquake}{earthquake} is a sudden shaking of the ground, caused by \omterm{def:g3:stones-pebbles-sand:rock}{rocks} that break
and slip deep underground. It lasts seconds (a minute or two for the very
largest), not hours.
\end{solution}

\begin{solution}{exo:g4:volcanoes-earthquakes-first-look:3}
Drop, cover, hold on.
\end{solution}

\begin{solution}{exo:g4:volcanoes-earthquakes-first-look:4}
Around the Pacific \omterm{def:g2:water-on-earth:ocean}{Ocean}.
\end{solution}

\begin{solution}{exo:g4:volcanoes-earthquakes-first-look:5}
(a) True: only instruments notice most of them. (b) False: \omterm{def:g4:volcanoes-earthquakes-first-look:eruption}{lava} cools and
hardens into new \omterm{def:g3:stones-pebbles-sand:rock}{rock}. (c) False: running while the ground shakes is
dangerous (you can fall, and glass and objects fall from buildings); drop,
cover and hold on.
\end{solution}

\begin{solution}{exo:g4:volcanoes-earthquakes-first-look:6}
It comes from deep below the \omterm{def:g4:volcanoes-earthquakes-first-look:volcano}{volcano}, where \omterm{def:g3:stones-pebbles-sand:rock}{rock} has melted. It rises up
through the middle of the \omterm{def:g4:volcanoes-earthquakes-first-look:volcano}{volcano} and comes out at the top.
\end{solution}

\begin{solution}{exo:g4:volcanoes-earthquakes-first-look:7}
(a) Chile (Peru, Ecuador or Mexico are also right): its \omterm{def:g4:volcanoes-earthquakes-first-look:volcano}{volcanoes} run along
its whole length, near the Pacific \omterm{def:g1:landscapes-around-us:water}{coast}. (b) Hawaii, in the middle of the
Pacific \omterm{def:g2:water-on-earth:ocean}{Ocean}, far from the shaded ring.
\end{solution}

\begin{solution}{exo:g4:volcanoes-earthquakes-first-look:8}
$141 \div 12$ is a little less than 12 (because $12 \times 12 = 144$): about
12 large \omterm{def:g4:volcanoes-earthquakes-first-look:earthquake}{earthquakes} a month.
\end{solution}

\begin{solution}{exo:g4:volcanoes-earthquakes-first-look:9}
Out of 20, a quarter is 5, so three quarters is $3 \times 5 = 15$: about 15
in the ring, and about $20 - 15 = 5$ outside it.
\end{solution}

\begin{solution}{exo:g4:volcanoes-earthquakes-first-look:10}
He should have dropped onto his hands and knees, covered his head and neck
under his table, and held on to it until the shaking stopped. A doorway
does not protect him from falling or flying objects, and running while the
ground shakes can make him fall.
\end{solution}

\begin{solution}{exo:g4:volcanoes-earthquakes-first-look:11}
Outside the ring: $1214 - 837 = 377$ \omterm{def:g4:volcanoes-earthquakes-first-look:volcano}{volcanoes}. A third of 1\,200 is 400,
so two thirds of 1\,200 is 800. The ring holds 837, a little more than 800
out of about 1\,200: ``about 2 out of 3'' is a good description.
\end{solution}
